
\documentclass[11pt,a4paper]{report}   % report: a document with chapters
\usepackage[utf8]{inputenc}            % allows letters with accents
\usepackage{graphicx}                  % for figures
\usepackage[margin=2.5cm]{geometry}    % page margins

\usepackage{titlesec}
\titleformat{\chapter}[hang]{\normalfont\huge\bfseries}{\thechapter.}{0.5em}{}
\titlespacing*{\chapter}{0pt}{0pt}{20pt}   % space: left, above, below the title

\usepackage[english]{babel}
\usepackage{csquotes}
\usepackage[style=apa, backend=biber]{biblatex}
\addbibresource{references.bib}

% Makes citations, URLs and DOIs clickable in the PDF (keep this package last).
% colorlinks gives coloured text instead of boxes; use [hidelinks] for plain black links.
\usepackage[colorlinks=true, citecolor=blue, linkcolor=black, urlcolor=blue]{hyperref}

\title{Point Taken: Can AI reduce annotation costs for Sample-based area estimation for landcover classification change}
\author{Niek Wels}
\date{\today}


\begin{document}

\maketitle   % prints the title page from \title, \author and \date above

\chapter{Introduction}

Deforestation is a significant cause of greenhouse gas emissions \parencite{masolele2021}. Monitoring
deforestation is mandated by international agreements, for participating countries \parencite{voigt2025}, such as the REDD+
(Reducing Emissions from Deforestation and Forest Degradation) programme of the United Nations Framework
Convention on Climate Change (UNFCCC). This programme requires developing countries to report changes in land-cover and offers
subsidies to countries that demonstrate a verified reduction in carbon emissions resulting from the land-cover change.

To verify this reduction, Sample-Based Area Estimation (SBAE) is used. SBAE is
a statistical method that estimates the area of land-cover classes from a
probability sample of random observations \parencite{jonckheere2024}. For
land-cover change, the area of change is inferred from the sample points
observed to change class between dates.

Pixel counting from a classified map gives biased area estimates, as omission
and commission errors do not cancel out \parencite{olofsson2014}. 
SBAE avoids this bias by estimating area from reference labels at randomly
selected sample points, each weighted by the area it represents \parencite{olofsson2014}.
Secondly, SBAE quantifies the uncertainty of the area estimate as a confidence
interval, which pixel counting cannot provide \parencite{olofsson2014}. Reporting
this uncertainty is required under the UNFCCC \parencite{ipcc2019}, and higher
uncertainty can reduce REDD+ results-based payments \parencite{jonckheere2024}.
%add more reasons we can think of during supervisor meeting this way it will directly be clear why this would work better than classifying

Because higher uncertainty can reduce REDD+ results-based payments, uncertainty
should be kept as low as possible. The most direct way to achieve this is to
increase the number of sample points. As land-cover change is rare, estimating
it precisely requires even more samples \parencite{pagliarella2018}.

Currently, sample points are labelled by human annotators using
high-resolution imagery such as Google Earth \parencite{olofsson2014}.
This is time-consuming, and for the large samples needed to estimate rare
change it is often not feasible within reporting deadlines
\parencite{jonckheere2024}. In this thesis, annotation costs therefore refer to the
time annotators spend labelling the sample points of an SBAE.
Machine-learning approaches, in particular deep learning, can reduce this
burden by labelling part of the sample automatically, so that a sufficient
number of sample points is reached in less time \parencite{jonckheere2024}.

Human interpretation is also subjective, as interpreters can assign different
labels to the same sample point \parencite{pengra2020, gfoi2018}. Such errors
can bias area estimates, and variability among interpreters increases their
uncertainty \parencite{masolele2021, jonckheere2024}. A model applies the
same criteria to every sample point.

One promising type of such models is the foundation model. Foundation models are
pretrained in a self-supervised way on large amounts of broad, unlabelled data,
through which they learn general patterns and relationships in that data
\parencite{szwarcman2025}. This thesis uses a geospatial foundation model: a
foundation model pretrained on large amounts of satellite imagery. It is based on
the Vision Transformer (ViT) architecture \parencite{dosovitskiy2021} and was
pretrained by reconstructing masked patches of satellite images
\parencite{szwarcman2025}. Because of these previously learned patterns, geospatial foundation models can
outperform deep learning models trained from scratch, such as CNNs
 \parencite{szwarcman2025,simumba2025}, and require
less labelled data to achieve accurate predictions \parencite{szwarcman2025}.

The labelled data can be used to fine-tune the encoder. Fine-tuning adapts
the pretrained encoder to the specific task and data, in this case sample point classification.
Fine-tuning slightly adjusts some of the encoder’s weights. This changes the embeddings it produces.
The embeddings are the output of the encoder, it creates the 
spatial patterns and relationships the model found in the data.
A classification head is added to the model architecture, since
the model is made for image reconstruction \parencite{szwarcman2025},
but needs to classify sample points. 
These adjusted embeddings help the classification head predict classes better \parencite{simumba2025}.

Even so, the fine-tuned model will not classify every sample point correctly, and
incorrect labels in SBAE can bias the area estimate \parencite{mcroberts2018}.
A human interpreter can indicate when they are unsure about a point; for a
model, this has to be derived from its output in the form of a confidence
score. In this thesis, confidence is expressed as the margin: the difference
between the probabilities of the two most likely classes \parencite{valle2023}.
A large margin indicates that the model clearly prefers one class, while a
small margin indicates doubt between the first and second class. It should
therefore be checked whether the confidence reflects the accuracy, meaning that
sample points with a higher confidence are more often classified correctly
\parencite{geifman2017}. If so, the confidence can be used to determine which
points still require human interpretation, reducing the annotation costs.





%from a probability sample of reference observations, 
%rather than by counting pixels in a classified map.”




% NOTES (from before, kept as comments so they do not show in the PDF):
% Copy the feedback of Susanne its very helpfull use that part of the introduction and critically look at what is
% needsed and whats not
% Notes: Focus more on Ai in the introduction
% Explain why SBAE is superior over landcover classification segmentation over the entire area.
% One reason is that REDD just requires condifdence intervals, which just landcover mapping wont give you, also due
% to the interest in rare classes (change), it could not work when doing landcover change since that is harder to
% model thus more errors
% by sending the confident AI points to a human it wont lose accuracy, this is important to state so also do not
% claim it will become more accuracte just more efficeitn


% Prints the reference list. Only cited sources appear;
% use \nocite{*} above this line to list every entry in references.bib.
\printbibliography

\end{document}
